\chapter{Board Bring-up}
\label{chap:fpgasteel-board-bringup}

\section{Introduction}

This chapter brings up the \term{DE1-SoC} board. It builds the bootloader and the \gterm{fpga}{FPGA} bitstream, and prepares the bootable \term{SD} card that the example programs of the following chapters are loaded from. The discussion begins with an overview of the board's boot chain, from reset to a running bare-metal application (Section~\ref{sec:fpgasteel-board-introduction}), then turns to the more practical side: preparing and building the bootloader itself. Unlike \prjname{FPGAlix}, no \term{Linux} kernel or root filesystem is involved: \gterm{uboot}{U-Boot} loads the compiled application directly, either as a \repopath{.bin} from the \term{SD} card at standalone boot or as an \repopath{.elf} pushed over \gterm{jtag}{JTAG} for debugging. The chapter also covers generating the \gterm{fpga}{FPGA} bitstream and the two boot scripts that select between these paths, before closing by assembling the bootloader, the bitstream and a boot script onto a bootable \term{SD} card, producing a working baseline ready for the examples developed in the remaining chapters.

\begin{table}[htbp]
\centering
\begin{tabularx}{\textwidth}{@{} l X @{}}
\toprule
\textbf{Component} & \textbf{Purpose} \\
\midrule
\gterm{uboot}{U-Boot} & Boots the board: loads the \gterm{fpga}{FPGA} bitstream, programs the fabric through the \gterm{fpgamanager}{FPGA Manager}\footnote{\glsentrydesc{fpgamanager}}, and loads the bare-metal application, either standalone from the \term{SD} card or over \gterm{jtag}{JTAG} for debugging. \\
\gterm{fpga}{FPGA} bitstream & Configures the \gterm{fpga}{FPGA} fabric with the custom hardware described in Chapter~\ref{chap:fpgasteel-hw-overview}. \\
Boot script & Selects which of the two boot paths above \gterm{uboot}{U-Boot} follows; see Table~\ref{tab:fpgasteel-boot-scripts}. \\
\bottomrule
\end{tabularx}
\caption{Main components built in this chapter, and their purpose.}
\label{tab:fpgasteel-board-bringup-components}
\end{table}

\section{Required software}
\label{sec:fpgasteel-board-required-software}

The following must be installed and available on the \term{PATH}:

\begin{table}[H]
\centering
\begin{tabularx}{\textwidth}{@{} l X @{}}
\toprule
\textbf{Software} & \textbf{Why it is needed} \\
\midrule
Ubuntu Desktop \term{x86\_64} & Host operating system the build steps in this chapter run on. \\
\gterm{quartusprime}{Quartus Prime} & Compiles the \gterm{fpga}{FPGA} hardware project and converts it into the bitstream loaded by \gterm{uboot}{U-Boot} (Section~\ref{sec:fpgasteel-generating-bitstream}). \\
\term{ARM} toolchain (xPack Arm Embedded GCC) & Cross-compiles \gterm{uboot}{U-Boot} and the bare-metal example applications for the \term{DE1-SoC}'s \term{ARM} \term{Cortex-A9} cores. \\
\bottomrule
\end{tabularx}
\caption{Required host software for this chapter, and why each is needed.}
\label{tab:fpgasteel-board-bringup-software}
\end{table}

See Paragraph~\ref{subsec:fpgasteel-install-paths} for versions and install paths.

A number of \gterm{apt}{apt}\footnote{\glsentrydesc{apt}} packages are also required, for building the bootloader and the example software. See Paragraph~\ref{subsec:fpgasteel-apt-packages} for the full list and Paragraph~\ref{subsec:fpgasteel-env-vars} for setting the \term{PATH}.

\section{Boot chain overview}
\label{sec:fpgasteel-board-introduction}

Even for a bare-metal application, the \gterm{hps}{HPS} cannot jump straight to user code out of reset: the \term{ARM} core has no working \term{DRAM}, clocks or pin muxing yet, and those can only be brought up by running code from on-chip \term{RAM} first. The bootloader (\gterm{spl}{SPL} + \gterm{uboot}{U-Boot}) performs this early hardware init and then loads the bare-metal application into \term{DDR}; it is not tied to running \term{Linux}.

On reset, the fixed-function \gterm{bootrom}{Boot ROM} samples the appropriate \gterm{fpga}{FPGA}'s pins to pick the boot source (\term{SD/MMC}, \term{QSPI}, \term{NAND}, or \gterm{fpga}{FPGA} fallback; see \cite[``Boot Select'']{CV-TRM}). On the \term{DE1-SoC} schematic \cite{DE1-SCH}, these pins are called \hw{HPS_BOOTSEL[2:0]} and are hardwired on the \gterm{pcb}{PCB} to 101 (boot from \term{SD/MMC}), so the \gterm{bootrom}{Boot ROM} reads the combined \gterm{spl}{SPL} (Secondary Program Loader) + \gterm{uboot}{U-Boot} image from a dedicated raw partition (type A2) on the \term{SD} card, loads it into on-chip \term{RAM}, verifies it, and hands off to \gterm{spl}{SPL}, which brings up clocks, pin muxing and \term{DDR}. \gterm{spl}{SPL} then loads full \gterm{uboot}{U-Boot} into \term{DDR}, which loads the \gterm{fpga}{FPGA} bitstream and programs the fabric through the \gterm{fpgamanager}{FPGA Manager}, then loads and starts the bare-metal application from wherever it is at that point: as an \repopath{.elf} pushed over \gterm{jtag}{JTAG} when debugging (e.g.\ from \term{Eclipse}/\gterm{openocd}{OpenOCD}, \gterm{uboot}{U-Boot} only needs to have brought up \term{DDR}), or as a raw \repopath{.bin} read from the \term{SD} card boot partition when booting standalone.

\begin{figure}[htbp]
    \centering
    \includesvg[width=\linewidth]{FPGAsteelBootChain}
    \caption{\term{HPS} boot chain from reset to the bare-metal application.}
    \label{fig:fpgasteel-boot-chain}
\end{figure}

\section{Bootloader generation flow}
\label{sec:fpgasteel-bootloader-generation}

\term{SoC EDS} (Embedded Design Suite), \gterm{altera}{Altera}'s\footnote{\glsentrydesc{altera}}, later \term{Intel}'s, companion toolset to \term{Quartus} for \term{SoC} \gterm{fpga}{FPGA} embedded development, is no longer available. Since \term{SoC EDS} v19.1 Std / v19.3 Pro, the bootloader (\gterm{uboot}{U-Boot}) build flow no longer relies on its \gterm{gui}{GUI}\footnote{\glsentrydesc{gui}} tool, \cmdpath{bsp-editor}, to configure and generate \gterm{uboot}{U-Boot}. The current flow is:

\begin{enumerate}
    \item Compile the hardware project in \gterm{quartusprime}{Quartus Prime}: this produces a handoff folder, a set of \term{XML} files describing how the \gterm{hps}{HPS} is configured (pin muxing, clocks, \term{SDRAM} timings, etc.).
    \item Run \cmd{cv_bsp_generator.py} \cite{MACNICA-BOOT} (part of the \gterm{uboot}{U-Boot} source tree) on the handoff folder: it generates the board-specific header files (\repopath{pinmux_config.h}, \repopath{pll_config.h}, \repopath{sdram_config.h}, \repopath{iocsr_config.h}).
    \item Copy these headers into \gterm{uboot}{U-Boot}'s generic source code, together with any custom user options (device tree, defconfig, etc.).
    \item Run \cmdpath{make} to build the bootloader image.
\end{enumerate}

\textbf{Note:} in other projects this flow normally starts by cloning the \gterm{uboot}{U-Boot} source from \gterm{github}{GitHub}. In \prjname{FPGAsteel} this step is not needed: \repopath{repos/u-boot-socfpga} \cite{SUB-UBOOT} is already tracked as a \gterm{git}{git}\footnote{\glsentrydesc{git}} submodule (see \repopath{.gitmodules}), so the source is already in place after cloning the repository.

\begin{warningbox}
    \gterm{platformdesigner}{Platform Designer}'s \term{IP} search path (Tools > Options > IP Catalog, ``IP Search Path'') is a \term{Quartus}-wide setting, not scoped to a single project. Every custom component this design uses (\hw{pclk_reset_controller}, \hw{ov7670_data_interface}, \hw{avalon_st_64_to_24_repacker}, \hw{pll_lock_reset_sync}, Hardware Overview~\ref{chap:fpgasteel-hw-overview}) lives under \repopath{hw/my_vhdl/qsys_components/}, and \gterm{platformdesigner}{Platform Designer} can only resolve them if that folder is listed there. Check this setting, not just the first time on a new machine, but every time you switch between working on \prjname{FPGAsteel} and \prjname{FPGAlix} on the same machine: each project has its own \repopath{hw/my_vhdl/qsys_components/}, and a search path left pointing at the other project's folder will make step 1 above fail.
\end{warningbox}

\begin{figure}[htbp]
    \centering
    \includesvg[width=\linewidth]{FPGAsteelBootloaderGeneration}
    \caption{Current bootloader generation flow (\term{SoC EDS} v19.1 Std / v19.3 Pro and later). Dark Gray boxes are \term{Quartus Prime}; green boxes are part of \term{U-Boot}.}
    \label{fig:fpgasteel-bootloader-generation}
\end{figure}

\section{Getting the repository}
\label{sec:fpgasteel-getting-repository}

The sections that follow assume the \prjname{FPGAsteel} repository is already cloned locally. Here's how, submodules included:

\begin{shellcodehost}
git clone --recurse-submodules https://github.com/tceccarini/FPGAsteel.git
cd FPGAsteel
\end{shellcodehost}

The option |--recurse-submodules| makes \cmdpath{git} also clone every submodule listed in \repopath{.gitmodules} right away --- these live under \repopath{repos/} (including |u-boot-socfpga|, needed in the next section). Without it, those folders would be created but left empty.

\section{Building the bootloader (U-Boot)}
\label{sec:fpgasteel-building-bootloader}

These are the concrete commands behind Figure~\ref{fig:fpgasteel-bootloader-generation}, run from the root of the repository.

\subsection{Generate the BSP headers from the Quartus handoff}
\label{subsec:fpgasteel-bsp-headers}

Compiling the hardware project makes \term{Quartus} export the \gterm{hps}{HPS} configuration as \term{XML} under the handoff folder placed at \repopath{hw/quartus/hps_isw_handoff/soc_system_hps_0/}. Convert this into the \term{C} headers \gterm{uboot}{U-Boot} needs by running:

\begin{shellcodehost}
python3 repos/u-boot-socfpga/arch/arm/mach-socfpga/cv_bsp_generator/cv_bsp_generator.py -i hw/quartus/hps_isw_handoff/soc_system_hps_0/ -o repos/u-boot-socfpga/board/terasic/de1-soc/qts/
\end{shellcodehost}

\subsection{Configure and build U-Boot}
\label{subsec:fpgasteel-configure-build-uboot}

\begin{shellcodehost}
cd repos/u-boot-socfpga
export ARCH=arm
export CROSS_COMPILE=arm-none-eabi-
make socfpga_cyclone5_defconfig
make menuconfig
\end{shellcodehost}

In \gterm{menuconfig}{menuconfig}\footnote{\glsentrydesc{menuconfig}}, select the board target: \term{ARM} architecture $\rightarrow$ \gterm{altera}{Altera} \term{SOCFPGA} board select $\rightarrow$ \term{Terasic} \term{DE1-SoC} (Cyclone V).

Exit and confirm \term{Yes} when prompted to save --- skipping this leaves the wrong board selected and the build silently targets the wrong hardware.

\begin{shellcodehost}
make -j$(nproc)
\end{shellcodehost}

\subsection{Result}
\label{subsec:fpgasteel-build-result}

The build produces \repopath{repos/u-boot-socfpga/u-boot-with-spl.sfp}, an image combining \gterm{spl}{SPL} and \gterm{uboot}{U-Boot} in the \gterm{altera}{Altera} \term{SoCFPGA} format. This will be the file written raw to the \term{SD} card's boot partition (see Figure~\ref{fig:fpgasteel-boot-chain}), covered in a dedicated section on preparing the \term{SD} card.

\section{Generating the FPGA bitstream (RBF)}
\label{sec:fpgasteel-generating-bitstream}

The \gterm{fpga}{FPGA} fabric on the \term{DE1-SoC} can be configured through several paths: over \gterm{jtag}{JTAG} from \term{Quartus}, over \gterm{jtag}{JTAG} via \gterm{openocd}{OpenOCD}, or by the \gterm{hps}{HPS} itself at boot through the \gterm{fpgamanager}{FPGA Manager}.

\prjname{FPGAsteel} uses the last one: during software development, \gterm{uboot}{U-Boot} loads a bitstream from the \term{SD} card and programs the fabric through the \gterm{fpgamanager}{FPGA Manager} before handing off to the application. Which path is used is a hardware setting, selected by \term{MSEL}.

\subsection{MSEL configuration (SW10)}
\label{subsec:fpgasteel-msel}

Using the \gterm{fpgamanager}{FPGA Manager} path, the one \gterm{uboot}{U-Boot} uses to program the \gterm{fpga}{FPGA}, requires setting \term{MSEL[4:0]} accordingly. On the \term{DE1-SoC} this is done with the \term{SW10} \term{DIP} switch, on the underside of the board.

\textbf{Note:} \term{SW10} uses negative logic: switch \term{ON} means logic 0, switch \term{OFF} means logic 1.

Set the six \term{SW10} switches, on the underside of the board, to the positions given in Table~\ref{tab:fpgasteel-sw10-msel}, which select the \gterm{fpgamanager}{FPGA Manager} path (\term{MSEL[4:0]} = 00000):

\begin{table}[htbp]
\centering
\begin{tabular}{@{} l l l l @{}}
\toprule
\textbf{Switch} & \textbf{Signal} & \textbf{Position} & \textbf{Logic} \\
\midrule
\term{SW10.1} & \term{MSEL0} & \term{ON} & 0 \\
\term{SW10.2} & \term{MSEL1} & \term{ON} & 0 \\
\term{SW10.3} & \term{MSEL2} & \term{ON} & 0 \\
\term{SW10.4} & \term{MSEL3} & \term{ON} & 0 \\
\term{SW10.5} & \term{MSEL4} & \term{ON} & 0 \\
\term{SW10.6} & --- & \term{OFF} & 1 \\
\bottomrule
\end{tabular}
\caption{\term{SW10} switch positions for the \term{FPGA Manager} path (\term{MSEL[4:0]} = 00000).}
\label{tab:fpgasteel-sw10-msel}
\end{table}

With \term{MSEL[4:0]} = 00000, the \gterm{fpga}{FPGA} is configured by the \gterm{hps}{HPS} via the \gterm{fpgamanager}{FPGA Manager}, as described above.

\subsection{Converting the SOF to RBF}
\label{subsec:fpgasteel-sof-to-rbf}

\term{Quartus} compiles the hardware project into a \repopath{.sof} (\term{SRAM} Object File), the full \gterm{fpga}{FPGA} configuration, in a format only \term{Quartus}/\gterm{jtag}{JTAG} tools understand. The \gterm{fpgamanager}{FPGA Manager} instead needs an \repopath{.rbf} (Raw Binary File), a stripped-down version of the same configuration. Convert it from the root of the repository:

\begin{shellcodehost}
quartus_cpf -c hw/quartus/output_files/FPGAsteel.sof sdcard/boot_partition/soc_system.rbf
\end{shellcodehost}

The output lands directly under \repopath{sdcard/boot_partition/}, next to the other files that will end up on the \term{SD} card's boot partition.

\section{Boot scripts}
\label{sec:fpgasteel-boot-scripts}

\gterm{uboot}{U-Boot} does not run a compiled \term{ELF} (or binary) directly at standalone boot: it runs a small boot script first, which tells it what to load and how to start it. \prjname{FPGAsteel} ships two of them under \repopath{sw/}, for different scenarios:

\begin{table}[htbp]
\centering
\begin{tabular}{@{} l p{0.72\linewidth} @{}}
\toprule
\textbf{Script} & \textbf{Purpose} \\
\midrule
\repopath{multiboot.script} & Standalone boot. Shows a menu of the built example apps, loads the \gterm{fpga}{FPGA} bitstream, loads the chosen app's \repopath{.bin} from the boot partition, and jumps to it with the \gterm{uboot}{U-Boot} command |go|. This is the ``Standalone: .bin'' path in Figure~\ref{fig:fpgasteel-boot-chain}. \\
\repopath{boot.script.jtag_debug} & One-time setup for \gterm{jtag}{JTAG} debugging. Programs the \gterm{fpga}{FPGA}, disables the \gterm{hps}{HPS}'s caches, and halts with a message, ready for \gterm{openocd}{OpenOCD}/\gterm{gdb}{GDB}\footnote{\glsentrydesc{gdb}} to load the \term{ELF}. This is the ``Debug: .elf'' path in Figure~\ref{fig:fpgasteel-boot-chain}. Used by every debug session, regardless of whether the example being debugged enables its own caches at runtime. \\
\bottomrule
\end{tabular}
\caption{\term{U-Boot} boot scripts shipped with \prjname{FPGAsteel}.}
\label{tab:fpgasteel-boot-scripts}
\end{table}

Both are self-installing: on first run they rewrite \gterm{uboot}{U-Boot}'s \sysid{bootcmd} environment variable and reboot, so from the second boot onward the board goes straight into that behavior without re-reading the script.

Only one script can be active at a time, whichever one was last compiled to \repopath{u-boot.scr}, the fixed filename \gterm{uboot}{U-Boot} looks for on the boot partition. Compile the one needed:

\begin{shellcodehost}
cd sw/
mkimage -A arm -O linux -T script -C none -a 0 -e 0 -n "Boot" -d multiboot.script ../sdcard/boot_partition/u-boot.scr
\end{shellcodehost}

Replace \repopath{multiboot.script} with \repopath{boot.script.jtag_debug} to switch boot mode. This step must be repeated every time you switch to a different script.

\section{Preparing the SD card}
\label{sec:fpgasteel-preparing-sd-card}

\begin{warningbox}
    This section involves low-level disk operations. A wrong device name or a mistyped command can silently overwrite a system disk, a data partition, or any other block device on your machine. Always double-check every device path (\extpath{/dev/sdX}, \extpath{/dev/loopN}) before pressing Enter. When in doubt, stop and verify.
\end{warningbox}

Unlike a \term{Linux} boot, a bare-metal \term{SD} card does not need a root filesystem partition: only the bootloader and the boot partition contents matter. All the files that go on the card are collected under \repopath{sdcard/}. Build a disk image file first, then write it to the physical \term{SD} card; the same procedure can be applied directly on the \term{SD} card by replacing \repopath{sdcard/images/sdcard.img} with the \term{SD} card device (e.g.\ \extpath{/dev/sdX}), skipping the image creation step and the final \cmdpath{dd}/\term{Etcher} write.

To write the final image to the \term{SD} card, use:

\begin{itemize}
    \item \term{Linux}: \cmdpath{dd} (see last step below)
    \item \term{Windows}/\term{Mac}: \gterm{balenaetcher}{balenaEtcher}\footnote{\gterm{balenaetcher}{balenaEtcher} --- \url{https://etcher.balena.io/}}, free and open-source, just select the image and the target drive
\end{itemize}

\subsection{Create the image file}
\label{subsec:fpgasteel-create-image}

From the root of the repository, move into the \repopath{sdcard/} directory:

\begin{shellcodehost}
cd sdcard/
mkdir -p images
\end{shellcodehost}

Create an empty image sized to fit any ``8 GB'' \term{SD} card. \term{SD} card manufacturers use decimal units (1 GB = 1{,}000{,}000{,}000 bytes), so an ``8 GB'' card has about 7629 MiB of actual space. Using 7400 MiB leaves a safe margin:

\begin{shellcodehost}
dd if=/dev/zero of=images/sdcard.img bs=1M count=7400
\end{shellcodehost}

\subsection{Partition the image}
\label{subsec:fpgasteel-partition-image}

Two partitions are needed: a \term{FAT32} boot partition (number 1) and a raw \term{A2} partition for \gterm{uboot}{U-Boot} (number 3). Partition number 2 is simply skipped; its space stays unallocated between the two.

\begin{shellcodehost}
fdisk images/sdcard.img
\end{shellcodehost}

Within \cmdpath{fdisk}, enter the following commands in order:

\begin{shellcodehost}
o       # new MBR partition table
n       # new partition 1, FAT32 boot (500 MiB)
p
1
(enter)
+500M
t       # change partition type
c       # set type to W95 FAT32 LBA (partition 1, FAT32 boot, 500 MiB)
p       # print disk info, note "Sectors" and "Sector size"
\end{shellcodehost}

Prior to creating partition 3, calculate its first sector:

\[
\text{first\_sector}_{p3} = \text{total\_sectors} - \left\lceil \frac{10 \times 1024 \times 1024}{\text{sector\_size}} \right\rceil
\]

Example, assuming 512-byte sectors and a 7400 \term{MiB} image (total sectors = 15{,}155{,}200):

\[
\begin{aligned}
\text{first\_sector}_{p3} &= 15{,}155{,}200 - \left\lceil \frac{10 \times 1024 \times 1024}{512} \right\rceil \\
&= 15{,}155{,}200 - 20{,}480 \\
&= 15{,}134{,}720
\end{aligned}
\]

Create partition 3 at that position:

\begin{shellcodehost}
n               # partition 3, raw A2, U-Boot SPL (10 MiB)
p
3
15134720        # enter the calculated first sector
(enter)         # accept default last sector (end of disk), fdisk should confirm ~10 MiB
t
3
a2              # set type to A2 (Altera SoCFPGA bootloader)
p               # verify the partition table before writing
w               # write and exit
\end{shellcodehost}

\subsection{Attach the image to a loop device}
\label{subsec:fpgasteel-attach-loop}

\begin{shellcodehost}
sudo losetup -fP images/sdcard.img
\end{shellcodehost}

Find out which loop device was assigned:

\begin{shellcodehost}
losetup -l
\end{shellcodehost}

Look for the entry with \repopath{sdcard.img} and note the device name (e.g.\ \extpath{/dev/loop2}). In all the commands below, replace \ccode{N} with the actual number.

\subsection{Format the partitions}
\label{subsec:fpgasteel-format-partitions}

\begin{shellcodehost}
sudo mkfs.vfat -n BOOT /dev/loopNp1
# loopNp3 is left unformatted, it will receive the raw U-Boot image with dd in the next step
\end{shellcodehost}

\subsection{Mount and copy the files}
\label{subsec:fpgasteel-mount-copy}

\begin{shellcodehost}
sudo mkdir -p /tmp/sdcard/boot
sudo mount /dev/loopNp1 /tmp/sdcard/boot
\end{shellcodehost}

Copy the boot partition files: the bitstream, the compiled boot script and the example binaries.

\begin{shellcodehost}
sudo cp -L sdcard/boot_partition/* /tmp/sdcard/boot/
\end{shellcodehost}

Write \gterm{uboot}{U-Boot} (\gterm{spl}{SPL} and \gterm{uboot}{U-Boot} proper, from Section~\ref{sec:fpgasteel-building-bootloader}) raw into the \term{A2} partition:

\begin{shellcodehost}
sudo dd if=sdcard/u-boot-with-spl.sfp of=/dev/loopNp3 bs=64k seek=0
sync
\end{shellcodehost}

\subsection{Unmount, detach and clean up}
\label{subsec:fpgasteel-unmount-detach}

\begin{shellcodehost}
sudo umount /tmp/sdcard/boot
sudo losetup -d /dev/loopN
sudo rm -rf /tmp/sdcard
\end{shellcodehost}

\subsection{Write the image to the SD card (Linux)}
\label{subsec:fpgasteel-write-sd}

\begin{shellcodehost}
sudo dd if=images/sdcard.img of=/dev/sdX bs=4M status=progress
sync
\end{shellcodehost}

Replace \extpath{/dev/sdX} with the actual \term{SD} card device. Double-check the target device before running: this command will overwrite it completely.
